% ===========================================================================
\section{Introduction}
\label{sec:introduction}
% ===========================================================================

National greenhouse-gas inventories provide the accounting basis for reporting under the
Paris Agreement. Under the Enhanced Transparency Framework, countries regularly report
sectoral emissions using Common Reporting Tables. In many developing countries,
agriculture, forestry and other land use (\AFOLU{}) accounts for a substantial share of
national emissions but is also among the sectors with the weakest activity data.

Pakistan illustrates this problem clearly. Agriculture is one of the country's largest
emitting sectors, yet its national inventory provides limited spatial detail. Some
livestock estimates also continue to rely partly on projections from older censuses.
National totals are necessary for international reporting, but they provide little
information about where emissions occur or where mitigation and improved data collection
should be targeted.

Agricultural activity also produces pollutants that affect air quality. Ammonia released
from fertiliser and manure is the principal alkaline precursor of secondary inorganic
particulate matter. Agricultural \NHthree{} can therefore contribute to PM$_{2.5}$
concentrations far from the fields and livestock systems where it originates.
Crop-residue burning in the Punjab rice--wheat belt also releases primary PM$_{2.5}$,
carbon monoxide and ozone precursors during the post-monsoon period, when particulate
pollution in Lahore is particularly severe.

Despite this importance, Pakistan has no gridded agricultural air-pollutant inventory
derived from national activity data. Air-quality modelling instead relies largely on global
emissions products. These products provide valuable regional coverage, but their activity
data and spatial assumptions cannot reliably distinguish, for example, an intensively
cultivated canal-belt district from an extensive rangeland district.

A useful national system must therefore meet three requirements. First, it must be
reproducible and updateable as new official statistics become available. Second, it must
provide provincial, district and gridded estimates that support subnational planning and
air-quality modelling. Third, its greenhouse-gas and air-pollutant calculations must use
consistent activity data and nitrogen flows.

The third requirement is physically important. A unit of manure nitrogen may be lost as
ammonia during housing, storage, grazing or field application, or it may remain in the
system and contribute to direct or indirect nitrous-oxide emissions. Calculating these
pathways in separate workbooks creates a risk of double counting, omission, or inconsistent
system boundaries. Similar inconsistencies arise when a national inventory, a
spatial-allocation exercise, and an emissions input for a chemistry model are developed
independently. The resulting products may appear detailed while failing to reproduce the
same national total or physical flow.

The same argument applies one level up. An agricultural inventory that is internally
consistent may still be inconsistent with the sectors around it, because several quantities
are genuinely shared: nitrogen leaving the manure-management balance as household fuel has to
arrive somewhere, and crop-residue burning has to be distinguishable from other open burning.
The framework described here is the agriculture module of a ten-sector national inventory
\citep{ahmad2026national}, in which agriculture is \SI{50.2}{\percent} of the national total
and those transfers are executable assertions rather than statements in a methods document.
This paper documents the sector; the companion documents what holds between sectors, and
neither repeats the other.

Many of the required inputs already exist. Pakistan publishes national and subnational
agricultural statistics, while open global datasets provide spatial information on
livestock, crops, land cover and fire occurrence. The remaining need is an integrated
framework that calculates emissions, allocates them spatially under an explicit
mass-balance constraint, evaluates the resulting products, and propagates uncertainty from
a common set of parameters.

This study develops such a framework for Pakistan. We construct a \emph{hybrid
Tier~1/Tier~2} agricultural inventory using the 2006 IPCC Guidelines and the 2019
Refinement. We then allocate individual emission streams to provinces, a
\SI{0.05}{\degree} grid, and administrative districts using process-specific spatial
proxies. Uncertainty is quantified using analytical propagation and Monte-Carlo simulation.

The central design principle is that the calculated inventory total serves as the control
value. Every provincial, gridded and district product is constrained to reconcile with that
control. Emission factors and activity data therefore determine the quantity emitted, while
spatial proxies determine only where those emissions are allocated.

\subsection{Contributions}
\label{sec:contributions}

This paper makes five main contributions.

First, it provides a reproducible national agricultural inventory implemented in a single
codebase that runs on the Python standard library alone. Emission factors, activity values
and other parameters are stored in versioned data files rather than embedded in the
calculation code. The inventory covers six IPCC categories: enteric fermentation (3.A),
manure management (3.B), rice cultivation (3.C), managed soils (3.D), field burning of
agricultural residues (3.F), and urea application (3.H). Liming emissions under category
3.G are not estimated: no national lime-use activity figure has been sought. All six
categories use Tier~1 equations with country-specific activity data. Selected livestock
parameters, including animal mass and manure-management-system shares, provide partial
Tier~2 characterisation. Tier status is therefore reported separately for each category.

Second, the framework allocates emissions to provinces, a national grid and administrative
districts under an explicit reconciliation constraint. Within each control area, every
emission stream is distributed in proportion to a process-matched proxy and then rescaled
so that the allocated values sum exactly to the control total. Livestock emissions are
allocated using gridded livestock density and, where available at district level, census
head-counts. Rice emissions are allocated using mapped paddy area, while residue-burning
emissions are allocated using satellite fire detections.

Third, the inventory is evaluated at several scales. National category totals are compared
with FAOSTAT, EDGAR and Climate TRACE using matched years and category definitions.
District patterns are evaluated against independent government crop-area records. The rice
allocation has a Spearman rank correlation of \num{0.92} with official district rice area,
while the managed-soils allocation reaches \num{0.84} against district wheat area.
Parametric uncertainty is estimated using both IPCC Approach~1 propagation and Monte-Carlo
aggregation. These methods apply different propagation formulations to the same underlying
uncertainty assumptions.

Fourth, the framework produces a national reconstruction for 1990--2024 and a
district-level historical panel based on fixed spatial shares. The national series reflects
changes in activity over time, while the district panel distributes each year's national
total according to the audited 2024 district pattern.

These are two different products on two different bases and are not interchangeable. The
canonical panel is the ten fiscal years FY2015--16 to FY2024--25, computed on Pakistani
fiscal activity with the reference year anchored to the 2024 census; it is what the release
reports, reconciles and prices uncertainty for. The 1990--2024 series is a contextual
reconstruction on FAOSTAT calendar-year activity, provided for long-run trend and for the
Common Reporting Table, and it is not the inventory's reporting basis. Matched-year differences separate the adopted stock/activity reconstruction from
the additional cotton-residue scope. Years lacking an archived calendar comparator
remain explicitly unpaired; no absent comparator is replaced by zero.

Fifth, the framework calculates agricultural air-pollutant co-emissions, including
\NHthree{}, \NOx{} and the pollutants released by crop-residue burning. These estimates use
the same activity data and process-specific spatial proxies as the greenhouse-gas
inventory. A shared, stage-resolved nitrogen ledger aligns the system boundaries used by
the IPCC and EMEP/EEA methods and helps diagnose differences between the resulting
estimates. The ammonia inventory is also compared with an independent gridded inventory for
South Asia. An exploratory Tier-1 land-sector extension is presented separately in
Supplementary Note~S1 and is excluded from all agricultural results.

Figure~\ref{fig:overview} summarises the framework. Section~\ref{sec:related_work} reviews
relevant inventory and spatial-allocation methods. Section~\ref{sec:data} describes the
study area and data. Sections~\ref{sec:methods}--\ref{sec:methods-shared} present the
greenhouse-gas inventory, spatial allocation, uncertainty analysis and shared nitrogen
accounting. Section~\ref{sec:validation} describes the evaluation framework.
Sections~\ref{sec:results} and~\ref{sec:airpollutants} report the greenhouse-gas and
air-pollutant results. Sections~\ref{sec:discussion}--\ref{sec:conclusion} discuss the
findings, limitations and conclusions.

\begin{figure*}[t]
\centering
\resizebox{0.96\textwidth}{!}{%
\begin{tikzpicture}[
  font=\footnotesize,
  band/.style={font=\scriptsize\bfseries\itshape, anchor=east, text=black!55},
  data/.style={draw,rounded corners,align=center,inner sep=5pt,fill=green!6,minimum height=11mm,
               text width=38mm},
  ledger/.style={draw,rounded corners,align=center,inner sep=6pt,fill=yellow!14,thick,
                 draw=black!55,text width=124mm,minimum height=11mm},
  calc/.style={draw,rounded corners,align=center,inner sep=5pt,minimum height=11mm,text width=52mm},
  ctrl/.style={draw,rounded corners,align=center,inner sep=5pt,fill=gray!10,minimum height=10mm,
               text width=62mm},
  prod/.style={draw,rounded corners,align=center,inner sep=5pt,fill=blue!7,minimum height=10mm,
              text width=32mm},
  cross/.style={draw,rounded corners,align=center,inner sep=5pt,fill=orange!8,text width=124mm,
                minimum height=8mm},
  arr/.style={-{Latex[length=2mm]},thick},
  ev/.style={-{Latex[length=2mm]},thick,dashed,black!55}
]
% ---- band 1: data, separated by function ----
\node[band] at (-3.6,7.4) {DATA};
\node[data] (act)  at (0,7.4)   {Inventory activity\\\scriptsize PBS census, NFDC,\\FAOSTAT history};
\node[data] (prox) at (4.3,7.4) {Spatial proxies\\\scriptsize GLW, WorldCover, FIRMS,\\MapSPAM (all held)};
\node[data] (eval) at (8.6,7.4) {Evaluation data\\\scriptsize MNFSR records, FAOSTAT/\\EDGAR/CT, \NHthree{} grid};

% ---- band 2: shared accounting ----
\node[band] at (-3.6,5.2) {SHARED ACCOUNTING};
\node[ledger] (led) at (4.3,5.2)
  {Livestock-N ledger \quad$|$\quad Fertiliser product mix \quad$|$\quad Crop-residue flows};

% ---- band 3: two calculations ----
\node[band] at (-3.6,3.2) {CALCULATION};
\node[calc,fill=blue!6]  (ghg) at (0.9,3.2) {Greenhouse gases\\\scriptsize IPCC 2006 + 2019 Refinement};
\node[calc,fill=red!6]   (air) at (7.7,3.2) {Air pollutants\\\scriptsize EMEP/EEA + IPCC};

% ---- band 4: controls and products ----
\node[band] at (-3.6,1.4) {CONTROLS};
\node[ctrl] (ctrl) at (4.3,1.4) {Stream-specific national totals\\$\rightarrow$ provincial controls};
\node[band] at (-3.6,-0.7) {PRODUCTS};
\node[prod] (grid) at (-0.4,-0.7) {\SI{0.05}{\degree} grid};
\node[prod] (dist) at (4.3,-0.7)  {ADM2 districts};
\node[prod] (ts)   at (9.0,-0.7)  {1990--2024 series};

% ---- cross-cutting ----
\node[band] at (-3.6,-2.6) {CROSS-CUTTING};
\node[cross] (cc) at (4.3,-2.6)
  {Mass balance \quad$|$\quad Parametric uncertainty \quad$|$\quad Structural scenarios};

\draw[arr] (act) -- (led);
\draw[arr] (led) -- (ghg);
\draw[arr] (led) -- (air);
\draw[arr] (ghg) -- (ctrl);
\draw[arr] (air) -- (ctrl);
\draw[arr] (ctrl) -- (grid);
\draw[arr] (ctrl) -- (dist);
\draw[arr] (ctrl) -- (ts);
% proxies drive the allocation, not the accounting: route down the left margin so the
% arrow does not cut through the shared-ledger or calculation boxes
\draw[arr] (prox.south) -- ++(0,-0.45) -| (-2.4,1.4) -- (ctrl.west);
\draw[arr,black!55,<->] (-3.05,-2.6) -- (-3.05,5.2)
  node[midway,rotate=90,anchor=south,font=\scriptsize\itshape,text=black!55]
  {structural scenarios feed back into the accounting};

% evaluation: short dashed arrows onto the products they evaluate
\draw[ev] (eval.east) -| (11.4,-0.7) -- (ts.east);
\draw[ev] (11.4,1.4) -- (ctrl.east);
\node[rotate=90,font=\scriptsize\itshape,text=black!55,anchor=south] at (11.75,3.4)
  {external evaluation};

\node[font=\scriptsize,align=center] at (4.3,-4.0)
  {$\sum_{p} E_{p,s} = E_{s}$ \qquad $\sum_{i\in p} e_{i,s} = E_{p,s}$
   \qquad $\sum_{d\in p} E_{d,s} = E_{p,s}$\\[1mm]
   \emph{grid and district products are parallel, not nested}};
\end{tikzpicture}%
}
\caption{%
Official activity data feed a shared livestock-nitrogen, fertiliser-nitrogen and crop-residue
accounting system. Separate greenhouse-gas (\CHfour{}, \NtwoO{}, \COtwo{}) and air-pollutant
(\NHthree{}, \NOx{}, PM$_{2.5}$, CO, NMVOC, BC, SO$_2$) calculations produce stream-specific
national totals, which become provincial controls and are allocated independently to a
\SI{0.05}{\degree} grid and to ADM2 districts on process-matched proxies; every product reconciles to
its control. Dashed arrows denote external evaluation, which is applied to the products rather than
to the calculations. Evaluation data are never used as inventory inputs.
}
\label{fig:overview}
\end{figure*}
